\hypertarget{result-and-scope}{%
\subsection{Result and scope}\label{result-and-scope}}

For every one of the 64 masks in
\nolinkurl{original_layout_dne_variants.json}, the numerical calculus
obtained by replacing the original fixed bundle with the indicated
normalized bundle has \textbf{exactly the same derivable numerical
codes} as the original calculus. This is a syntactic interderivability
theorem, uniform in all natural-number parameters. It includes the cases
in which a quantifier binds variable 0, the designated falsity code is
also a relation code, parameters collide, or the implication constructor
has the numerical alias \texttt{Imp(0,0)=0}.

The machine source must implement the corresponding new bundle
constructor. Its AX input reads, GEN row, eight logical rows, output
positions, gate and cleanup remain unchanged. The proof below supplies
the new logical/enumerator bridge. Combining it with a complete compiler
simulation and the exact blank-start frame for a resulting literal
binary table gives the same ZFC consistency equivalence as the original
source. An old/new source bisimulation that demands identical numerical
registers is neither needed nor generally true.

No Lean build of this new transfer theorem is supplied here. The
mathematical proof is explicit, and \texttt{check\_dne\_variants.py}
independently checks the finite template/emission facts for all 64
masks. Complete binary coverage and bootstrap verification belong to
each actual compiled candidate and must be reported separately.

ProveIt source basis: commit
\nolinkurl{6446ff5668307a2ac6aee1896037f0b792090389}. The exact formula
specification is the attributed upstream
\nolinkurl{ZF/m254/tools/logic/spec.py} at commit
\nolinkurl{297ab8df7e5de329605c3875b4934643b371a5e3}, SHA-256
\nolinkurl{7d96deff6994315b714579c74d10efeb65374694d476c802c727356e258578e9}.

\hypertarget{work-with-the-actual-numerical-calculus}{%
\subsection{1. Work with the actual numerical
calculus}\label{work-with-the-actual-numerical-calculus}}

Put

\[
P(x,y)=\frac{(x+y)(x+y+1)}2+y,
\quad I(x,y)=4P(x,y),\quad R(x,y)=4P(x,y)+1,
\] \[
Q_i A=4P(i,A)+2,\qquad E(x,y)=4P(x,y)+3,\qquad F=1.
\]

Use the abbreviations

\[
NA=I(A,F),\quad A\mathbin{\wedge_F}B=N I(A,NB),
\quad D_iA=NQ_iNA.
\]

These are functions on natural numbers. We do not assume they freely
generate a syntax tree. In particular, \texttt{I(0,0)=0} and
\texttt{R(0,0)=F} are genuine identities. Symbolic trees in the artifact
describe expressions formed from these functions; every identity of
derivability below survives arbitrary evaluation of their leaves.

Write \texttt{T\ ⊢\ A} for a theorem predicate closed under the
following common rules and rows:

\begin{itemize}
\tightlist
\item
  Modus ponens: from \texttt{A} and \texttt{I(A,B)}, infer \texttt{B}.
\item
  Generalization: from \texttt{A}, infer \texttt{Q\_i\ A}, for every
  natural \texttt{i}.
\item
  The Łukasiewicz row \[
  \rho_8(a,b,c,u)=((c\to u)\to a)\to((a\to c)\to(b\to c)).
  \]
\item
  The original falsity/quantifier row \[
  \rho_6(i,A)=(D_iQ_iA)\to A.
  \]
\item
  Quantifier distribution \[
  \rho_2(i,A,B)=Q_i(A\to B)\to(Q_iA\to Q_iB).
  \]
\end{itemize}

Both the original and every normalized calculus contain this common
core, without using either bundle as a premise.

\hypertarget{lemma-1-full-classical-implication-in-the-common-core}{%
\subsubsection{Lemma 1: full classical implication in the common
core}\label{lemma-1-full-classical-implication-in-the-common-core}}

The checked Łukasiewicz derivation supplies, uniformly in its arguments,

\[
K(A,B)=A\to(B\to A),
\] \[
S(A,B,C)=(A\to(B\to C))\to((A\to B)\to(A\to C)),
\] \[
\operatorname{Peirce}(A,B)=((A\to B)\to A)\to A.
\]

Only MP and instances of \texttt{rho8} occur in those derivations. The
repository states this polymorphically for an arbitrary binary
operation; injectivity of implication is not an assumption. See
\texttt{Lukasiewicz.lean}, the theorems \texttt{k}, \texttt{s}, and
\texttt{peirce}, and \texttt{Code254.lean}, \texttt{implicational} and
\texttt{c254\_classical}.

The common core also proves \texttt{F\ →\ A} for every natural
\texttt{A}. To see this directly, fix \texttt{i=0} and abbreviate

\[
H=Q_0NQ_0A,\qquad X=D_0Q_0A=H\to F.
\]

An instance of K gives \texttt{F\ →\ (H\ →\ F)}, hence \texttt{F\ →\ X}.
The row \texttt{rho6(0,A)} gives \texttt{X\ →\ A}. Transitivity of
implication, derived from K and S, gives \texttt{F\ →\ A}. This is
exactly the \texttt{Code254.c254\_efq} proof; it uses neither IR, atomic
vacuity, the CP/GIC bundle, nor a freshness condition.

\hypertarget{lemma-2-both-double-negation-directions}{%
\subsubsection{Lemma 2: both double-negation
directions}\label{lemma-2-both-double-negation-directions}}

For every natural \texttt{A}, the common core proves

\[
A\to NNA\qquad\hbox{and}\qquad NNA\to A.
\]

For the first direction, under assumptions \texttt{A} and \texttt{A→F},
MP yields \texttt{F}. Apply the propositional deduction theorem twice.

For the second, assume \texttt{NNA}. Under a further assumption
\texttt{A→F}, MP yields \texttt{F}, and \texttt{F→A} yields \texttt{A}.
Discharging the latter assumption gives \texttt{(A→F)→A}; Peirce gives
\texttt{A}. Discharge \texttt{NNA}. These uses of the deduction theorem
involve only MP and common-core theorems. Generalization is never
applied to an undischarged assumption.

\hypertarget{lemma-3-replacement-in-formula-contexts}{%
\subsubsection{Lemma 3: replacement in formula
contexts}\label{lemma-3-replacement-in-formula-contexts}}

Let \texttt{A\ ≃\ B} mean that both \texttt{A→B} and \texttt{B→A} are
common-core theorems. Then \texttt{≃} is an equivalence relation and is
preserved by the following context constructors:

\[
C[\,]\ ::=\ [\,]\mid I(C[\,],U)\mid I(U,C[\,])\mid Q_i C[\,].
\]

Implication congruence is propositional. For example, from
\texttt{A\textquotesingle{}→A} and \texttt{B→B\textquotesingle{}},
assumptions \texttt{A→B} and \texttt{A\textquotesingle{}} produce
\texttt{B\textquotesingle{}}; discharging gives
\texttt{(A→B)→(A\textquotesingle{}→B\textquotesingle{})}. Reverse the
two equivalences for the other direction.

For quantifier congruence, generalize the theorem \texttt{A→B} to obtain
\texttt{Q\_i(A→B)} and apply \texttt{rho2(i,A,B)}. This gives
\texttt{Q\_iA→Q\_iB}. Repeat for the converse. There is no condition on
\texttt{i}: this generalizes a theorem, including when \texttt{i=0} and
the literal number \texttt{F=1} decodes as \texttt{R(0,0)}.

No selected rewrite occurs in a variable-label argument of \texttt{Q},
\texttt{R}, or \texttt{E}. Thus these context constructors cover every
path used by the finite transformation.

\hypertarget{why-a-purely-semantic-shortcut-needs-care}{%
\subsubsection{Why a purely semantic shortcut needs
care}\label{why-a-purely-semantic-shortcut-needs-care}}

For an arbitrary propositional atom \texttt{f}, \texttt{((A→f)→f)} is
equivalent to \texttt{A∨f}, rather than to \texttt{A}. Five selected
redexes have an inner formula \texttt{A=Y→(Z→F)}; there \texttt{F→A} is
a tautology even when F is arbitrary, so local semantic cancellation
works. The remaining redex has a quantified inner formula. Moving from a
free assumption \texttt{¬R(0,0)} through a binder of variable 0 would be
invalid without further argument. Lemmas 1--3 avoid that issue
completely: they derive the required equivalences in a
bundle-independent numerical calculus and then apply GEN to proved
implications. The full six-redex theorem is consequently stronger and
cleaner than a claim restricted to the five locally absorbing redexes.

\hypertarget{the-exact-source-transformation}{%
\subsection{2. The exact source
transformation}\label{the-exact-source-transformation}}

Let \texttt{beta(a,b,c)} denote the original expression

\[
\bigl(C(a,b,c)\wedge_F G(a,b,c)\bigr)
\wedge_F
\bigl((R(a,b)\to Q_cR(a,b))\wedge_F(E(a,b)\to Q_cE(a,b))\bigr).
\]

The exact C and G expressions are taken directly from the pinned
\texttt{spec.py}, not reconstructed from an informal set-theory
description. Here \texttt{c} remains both the supplied GIC matrix code
and a variable-label parameter; its occurrences are not renamed or
reinterpreted.

An address is a list of child numbers in the fully expanded constructor
expression: child 1 is the first argument and child 2 is the second. The
six original addresses are:

\begin{longtable}[]{@{}ll@{}}
\toprule\noalign{}
Redex ID & Original address \\
\midrule\noalign{}
\endhead
\bottomrule\noalign{}
\endlastfoot
0 & \texttt{1,1,1,1,1,2} \\
1 & \texttt{1,1,1,1,1,2,1,1,1,2,2,1,1,1,2} \\
2 & \texttt{1,1,1,2} \\
3 & \texttt{1,1,1,2,1,1,2} \\
4 & \texttt{1,1,1,2,1,1,2,1,1,1,2,2,2,2,1,2} \\
5 & \texttt{1,2} \\
\end{longtable}

At a selected address replace \texttt{I(I(X,F),F)} by \texttt{X}.
Addresses refer to the original expression throughout, so deleting an
enclosing selected node does not accidentally select a different
descendant. Mask \texttt{m} selects ID \texttt{j} when bit \texttt{j} of
\texttt{m} is one. The recursive implementation transports the original
addresses through a deletion.

Write \texttt{beta\_m} for the result. Lemmas 2--3, by recursion on the
fixed expression, give

\[
\vdash_{\rm core}\beta(a,b,c)\to\beta_m(a,b,c),
\quad
\vdash_{\rm core}\beta_m(a,b,c)\to\beta(a,b,c)
\tag{1}
\]

for every mask and every \texttt{a,b,c}. All parameter equalities are
allowed in (1). The actual bundle-axiom gate remains
\texttt{c≠a\ ∧\ c≠b}.

Define \texttt{C\_m} to have exactly the original truth rule, MP, GEN
and eight rows, and the sole changed axiom rule

\[
c\ne a\ \land\ c\ne b\quad\Longrightarrow\quad C_m(\beta_m(a,b,c)).
\]

\hypertarget{theorem-4-equality-of-the-numerical-theorem-sets}{%
\subsubsection{Theorem 4: equality of the numerical theorem
sets}\label{theorem-4-equality-of-the-numerical-theorem-sets}}

For every natural \texttt{n} and every mask \texttt{m},

\[
C_m(n)\ \longleftrightarrow\ C_{254}(n).
\tag{2}
\]

Induct on an original derivation to obtain the forward translation into
the modified calculus. The truth, MP, GEN and eight-row cases are
unchanged. In a bundle case, use the new bundle axiom with the same
parameters and unchanged gate, followed by the common-core implication
\texttt{beta\_m→beta} from (1). The reverse translation is identical
with the implication reversed. This replaces each changed axiom leaf by
a finite derivation and preserves the conclusion code exactly.

This proves more than equiconsistency. In particular, every original
CP/GIC/vacuity projection, every original native-formula adapter, and
every original derivation of IR remains available in the new calculus,
with the same numerical conclusion. There is no need to reprove
set-theoretic calibration for a weakened schema or a changed variable
convention: (2) returns directly to the original literal calculus.

\hypertarget{exact-ax-interface-and-source-implementation}{%
\subsection{3. Exact AX interface and source
implementation}\label{exact-ax-interface-and-source-implementation}}

Use the original six boundary registers \texttt{(pl,np,t,S,d,h)}, where
the scratch registers and axiom selector have been cleared. Every
command begins with the unchanged fetch/read prefix. If \texttt{pl=0},
it selects candidate \texttt{np} and increments \texttt{np}; otherwise
it decrements \texttt{pl} before unpairing. It decodes the selector and
an additive field \texttt{x}, giving \texttt{d\textquotesingle{}=d+x}.

The three non-AX commands are unchanged:

\begin{itemize}
\tightlist
\item
  MP checks \texttt{I(t,d\textquotesingle{})=h}. Success sets
  \texttt{t=d\textquotesingle{}} and pushes zero onto \texttt{S};
  failure sets \texttt{t=0}. Either branch clears \texttt{h}.
\item
  DROP pushes the old \texttt{t} and sets \texttt{t=h}, retaining
  \texttt{h}.
\item
  NIP halts exactly when the incoming \texttt{h=1}; otherwise it pops
  \texttt{(h,S)=unP(S)}.
\end{itemize}

For AX, the selector supplies \texttt{c}; put
\texttt{a=d\textquotesingle{}}. Read the next field \texttt{y}, put
\texttt{b=h+y}, and read the same fourth field \texttt{u} as before.
Above the protected old stack tail, AX pushes, from bottom to top,

\[
t,\ Q_a h,\ \rho_8(a,b,c,u),\ \rho_7(a,b,c),\ \rho_6(a,b),\ \rho_5(a,b,c),
\] \[
\rho_4(a,b,c),\ \rho_3(a,b,c),\ \rho_2(a,b,c),\ \rho_1(a,b).
\]

The new current formula is \texttt{beta\_m(a,b,c)} if \texttt{c≠a,b},
and zero otherwise. Clear \texttt{d,h} and all scratch/selector
registers. \textbf{Only the displayed current bundle value changes.}

The unchanged AX word prefix is \texttt{{[}166,121,180,176{]}}: push
\texttt{a}, push the old cache, build its GEN row, then read the new
\texttt{b}. The eight original rows follow in original push order. The
only destructive \texttt{u} read remains inside the unchanged rho8
constructor. After the changed bundle, the original tail
\texttt{{[}192,131,211,197,131{]}} applies the same gate and cleanup.

The independent symbolic checker establishes these facts for every mask:

\begin{enumerate}
\def\labelenumi{\arabic{enumi}.}
\tightlist
\item
  The GEN prefix, eight complete row expressions and their order are
  literal matches to the pinned source.
\item
  The input reads are identical and occur at the same stack depths.
\item
  RPN execution never consumes the protected old current formula or the
  old stack tail.
\item
  The new constructor word leaves exactly one current formula above the
  same ten pushed outputs.
\item
  The bundle has outer constructor \texttt{I}, so its numerical code is
  divisible by four and is never \texttt{F=1}. Every pushed logical row
  also has outer \texttt{I}, except GEN, which has tag two; no pushed
  row is 1.
\item
  The gate and cleanup macros are byte-for-byte unchanged.
\end{enumerate}

To lift this symbolic identity to a changed structured register program,
apply the already proved scalar-push and binary-constructor contracts by
induction on the emitted RPN word, followed by the unchanged
gate/cleanup contracts. Each complete formula constructor is balanced:
it leaves its value as current, restores its scratch and parameter
registers, and pushes the previously current formula onto the protected
tail. A binary constructor combines exactly the two fresh arguments; no
new register operation is introduced. This is a compositional source
lemma for the new program, not a comparison asserting that it emits the
old numerical bundle.

The routed baseline's fine-token AST is not literally the old
\texttt{Dag254.v} AST: previous optimizations commute independent
increments and reassociate control nodes with corresponding continuation
adjustments. A separate all-register macro bridge now checks \textbf{all
19 entries} of \texttt{macro\_to\_fine.json}. It independently builds
finite control-flow graphs for the old macro and the composed fine word,
compares the nine-register increment vectors between matching decrement
guards, explores both branches, and verifies the resulting path
equations again with the upstream path decoder. The certificate has
\textbf{192 state pairs and 346 checked edges}. It preserves exact
output registers and return/divergence behavior. See the neighboring
machine-reduction artifacts \nolinkurl{audit_macro_fine_bridge.py},
\texttt{macro-fine-bridge.json}, and
\nolinkurl{macro-fine-bridge-certificates.json}. Composing that bridge with
the original macro contracts is the required justification for assigning
original constructor meanings to the new fine words; the aggregate
original AX word hash alone would not establish it.

The finite fine-token representation has these exact sizes:

\begin{longtable}[]{@{}
  >{\raggedright\arraybackslash}p{(\columnwidth - 8\tabcolsep) * \real{0.1579}}
  >{\raggedleft\arraybackslash}p{(\columnwidth - 8\tabcolsep) * \real{0.2105}}
  >{\raggedleft\arraybackslash}p{(\columnwidth - 8\tabcolsep) * \real{0.2105}}
  >{\raggedleft\arraybackslash}p{(\columnwidth - 8\tabcolsep) * \real{0.2105}}
  >{\raggedleft\arraybackslash}p{(\columnwidth - 8\tabcolsep) * \real{0.2105}}@{}}
\toprule\noalign{}
\begin{minipage}[b]{\linewidth}\raggedright
Mask
\end{minipage} & \begin{minipage}[b]{\linewidth}\raggedleft
Deleted redexes
\end{minipage} & \begin{minipage}[b]{\linewidth}\raggedleft
Constructor-expression nodes
\end{minipage} & \begin{minipage}[b]{\linewidth}\raggedleft
Wade macro calls
\end{minipage} & \begin{minipage}[b]{\linewidth}\raggedleft
Fine AX terminals
\end{minipage} \\
\midrule\noalign{}
\endhead
\bottomrule\noalign{}
\endlastfoot
0 & 0 & 145 & 186 & 1,270 \\
59 & 5, omitting ID 2 & 125 & 196 & 1,180 \\
63 & 6 & 121 & 196 & 1,162 \\
\end{longtable}

Every deleted redex removes 18 fine terminals and four
constructor-expression nodes. It can increase the count of coarser Wade
macro calls because an \texttt{And} or \texttt{Ex} pattern ceases to be
recognizable as a single derived constructor. Consequently source size
alone does not predict compressed state count.

\hypertarget{soundness-and-completeness-of-the-changed-enumerator}{%
\subsection{4. Soundness and completeness of the changed
enumerator}\label{soundness-and-completeness-of-the-changed-enumerator}}

\hypertarget{soundness-from-a-safe-start}{%
\subsubsection{Soundness from a safe
start}\label{soundness-from-a-safe-start}}

Let \texttt{Good\_m(S)} mean that every element obtained by repeatedly
unpairing the logical stack is a theorem of \texttt{C\_m}; the zero tail
contributes the truth code 0. Preserve

\[
C_m(t)\ \land\ C_m(h)\ \land\ Good_m(S).
\tag{3}
\]

The ordinary command proof is unchanged. Successful MP follows by MP in
\texttt{C\_m}; failed MP and pushed dummy zeros use truth. DROP and NIP
move previously proved codes. AX pushes GEN of the proved cache and
eight unchanged logical axiom instances, and its current formula is
either the gated new axiom or truth. Thus a halt, which is precisely an
incoming cache value \texttt{F}, implies \texttt{C\_m(F)}.

By (2), the sets of safe formula/stack contents for the original and new
calculi are equal. At the actual baseline seed
\texttt{(N,T,W,A,D,H,X,Y,L)=(3,0,0,0,3,0,0,0,0)}, all three
formula-bearing fields are zero, so (3) holds. The nonzero candidate
counter and additive parameter are not soundness assumptions about
formulas.

\hypertarget{completeness-by-explicit-proof-blocks}{%
\subsubsection{Completeness by explicit proof
blocks}\label{completeness-by-explicit-proof-blocks}}

Use the original contract: a command list \texttt{Build(A)} started from
\texttt{(t,S,d,h)=(T,S,0,0)} either halts, or ends at
\texttt{(A,S,0,0)}, for arbitrary \texttt{T,S}. Earlier halting is an
allowed successful outcome of this contract.

The unchanged stack row count and the tag facts above permit precisely
the original proof-block construction. In the following notation AX
\texttt{(c;x,y,u)} denotes selector \texttt{c+3} followed by the three
data fields; every unprinted command field is zero.

\begin{longtable}[]{@{}
  >{\raggedright\arraybackslash}p{(\columnwidth - 2\tabcolsep) * \real{0.5000}}
  >{\raggedright\arraybackslash}p{(\columnwidth - 2\tabcolsep) * \real{0.5000}}@{}}
\toprule\noalign{}
\begin{minipage}[b]{\linewidth}\raggedright
Derivation case
\end{minipage} & \begin{minipage}[b]{\linewidth}\raggedright
Command block
\end{minipage} \\
\midrule\noalign{}
\endhead
\bottomrule\noalign{}
\endlastfoot
Discard / truth & \texttt{MP(1);\ AX(1;0,0,0);\ NIP\^{}10} \\
Gated new bundle & \texttt{Discard;\ AX(c;a,b,0);\ NIP\^{}10} \\
Logical row \texttt{rho\_j}, \texttt{1≤j≤8} &
\texttt{Discard;\ AX(c;a,b,u);\ NIP\^{}j;\ DROP;\ NIP\^{}(11-j)} \\
GEN \texttt{i} from \texttt{A} &
\texttt{Build(A);\ DROP;\ NIP;\ AX(i;i,0,0);\ NIP\^{}9;\ DROP;\ NIP;\ NIP} \\
MP to \texttt{B} from \texttt{A} and \texttt{A→B} &
\texttt{Build(A);\ DROP;\ Build(A→B);\ NIP;\ DROP;\ NIP;\ MP(B);\ AX(B;0,0,0);\ NIP\^{}10;\ DROP;\ NIP;\ NIP} \\
\end{longtable}

The cleanup AX calls have \texttt{c=a}, so their bundle gate fails
exactly as before. In the row selector the new current bundle may be
pushed temporarily by DROP, but it is never 1, so the same NIP sequence
restores the protected tail. The bundle-leaf case uses the new
expression directly. Induction on a \texttt{C\_m} derivation therefore
gives the required block for every theorem code, including \texttt{F}.

\hypertarget{no-refutation-is-lost-by-beginning-with-nextproof-3}{%
\subsubsection{No refutation is lost by beginning with nextproof
3}\label{no-refutation-is-lost-by-beginning-with-nextproof-3}}

The precise original safe-start proof is in
\texttt{EnumSafeStart254.lean}, theorems \texttt{pad\_code\_ge},
\texttt{padded\_program\_halts}, \texttt{reach\_ge}, and
\texttt{complete\_from}. It adapts without altering any command
schedule.

Let \texttt{p} be a contradiction block. For every natural \texttt{u},
prepend

\[
\operatorname{AX}(0;1,1,u)
\]

and append \texttt{DROP;NIP;NIP}. The leading AX resets \texttt{d=h=0}
from arbitrary working values. Its current output and its newly pushed
stack may depend on those old values and on \texttt{u}, but the contract
for \texttt{p} quantifies over both. Thus this padded program halts from
every working state.

Its candidate code is at least \texttt{u}. Explicitly, if the encoded
remaining command list is \texttt{r}, the leading code is

\[
P(3,P(1,P(1,P(u,r)))),
\]

with the usual successor convention included in \texttt{r}; every Cantor
pairing is at least either argument. Pick \texttt{u} at least the
initial next-candidate counter. Unchanged decrement-before-unpairing
drains each current proof list, and every completed candidate increments
that counter. Hence the padded refutation is eventually selected unless
an earlier genuine halt occurs.

Consequently, from every start satisfying (3), and in particular the
actual seed with nextproof 3,

\[
\operatorname{Halt}(\operatorname{Enum}_m)\quad\longleftrightarrow\quad C_m(F)
\quad\longleftrightarrow\quad C_{254}(F).
\tag{4}
\]

This argument does not assume that the earlier candidates 0, 1, or 2 are
harmless, or that the current proof-list pointer is initially zero. It
supplies the needed cofinality directly.

\hypertarget{the-exact-cpgic-and-zfc-comparison}{%
\subsection{5. The exact CP/GIC and ZFC
comparison}\label{the-exact-cpgic-and-zfc-comparison}}

The inherited native theory has literal sealed sentences for the raw
container axiom CP, global-image collection GIC for \textbf{every actual
finite native formula with all parameters}, and raw irreflexivity IR.
The upstream source names these compact axioms C and G; the source
bundle also supplies atomic vacuity for the numerical calculus.

In ordinary set-theoretic notation their content is

\[
\forall p\,\exists V\Bigl(p\in V\ \land\
 \forall x\in V\,[x\subseteq V\ \land\ \exists W\in V\,\forall z\,(z\subseteq x\to z\in W)]\Bigr),
\]

and, for every formula \texttt{phi(x,y,parameters)},

\[
\forall B\,\exists C\Bigl[
 \forall x\in B\,\forall y\,(\phi(x,y)\to\exists y'\,(\phi(x,y')\land y'\in C))
 \ \land\ \forall y\in C\,\exists x\,\phi(x,y)\Bigr].
\]

IR is \texttt{∀x\ ¬(x∈x)}. These expressions describe the already
calibrated theory; the precise finite syntax and closures remain those
of \texttt{NativeTheory254.NativeAx}.

The proved numerical/native interface is

\[
C_{254}(F)\quad\longleftrightarrow\quad
\operatorname{BProv}(\operatorname{NativeAx},[\,],\bot).
\tag{5}
\]

Equation (2) gives exactly the same interface for \texttt{C\_m}, not
just a semantic soundness implication. All formula families, gate
conventions, equality transport, and falsity aliases are inherited
through the equality of theorem sets.

The ordinary ZFC calibration may be composed as in the baseline proof:
finite syntax and first-order completeness produce nonempty set-coded
models; the native CP/GIC model has the relevant bisimulation quotient
and internally well-founded interpretation; ZF including Foundation is
compared with ZFC by the internal constructible interpretation. Internal
Foundation must not be replaced by external well-foundedness of an
arbitrary countable model. The native theory is also not obtained merely
by deleting Choice from a displayed axiom list. Those are substantive
inherited interpretation arguments, referenced below, and are unchanged
by this transfer.

Within ZFC these are assertions about sets of finite proofs,
satisfaction for set structures, and finite machine runs. Equations
(1)--(4) are uniform finite derivation transformations and elementary
inductions and can be carried out there. Combined with the established
ordinary baseline interpretation, they yield

\[
\mathrm{ZFC}\vdash
\bigl(C_m(F)\leftrightarrow\neg\operatorname{Con}(\mathrm{ZFC})\bigr).
\tag{6}
\]

This is not an assertion that the new theorem has already been checked
in Lean, nor that every universe-polymorphic Lean result automatically
becomes a ZFC proof. The particular set constructions and inductions in
the ordinary baseline calibration, followed by the explicit finite
transformations here, are the mathematical reason for (6).

\hypertarget{what-is-needed-for-an-actual-binary-machine}{%
\subsection{6. What is needed for an actual binary
machine}\label{what-is-needed-for-an-actual-binary-machine}}

For a literal compiled candidate \texttt{M\_m}, the remaining
implementation theorem has the following concrete premises:

\begin{enumerate}
\def\labelenumi{\arabic{enumi}.}
\tightlist
\item
  The exact table has a complete positive-time simulation of every
  instruction of the \textbf{new} structured source, for arbitrary
  natural register values, all legal continuations, and the full tape
  frame. Both directions preserve genuine halting; neither a decode
  failure nor internal divergence is a halt.
\item
  The constructor/source lemma in Section 3 identifies that source with
  \texttt{Enum\_m}, and unchanged control fragments preserve the
  original fetch, selector, reads and command semantics.
\item
  The exact blank trace reaches the documented root frame, with the
  actual head, header, logical register values, pointer stack, exterior
  tape, and physical spare suffix. In particular the formula-bearing
  logical fields satisfy (3).
\end{enumerate}

These premises imply by simulation, (4), and (6)

\[
\mathrm{ZFC}\vdash
\bigl(\neg\operatorname{Halt}_{\rm blank}(M_m)
\leftrightarrow\operatorname{Con}(\mathrm{ZFC})\bigr).
\]

A complete instruction/compiler certificate plus bootstrap is sufficient
\textbf{when it is a certificate for the new source and is composed with
the new AX/enumerator lemma}. An audit of the old source alone does not
identify changed AX outputs. No equality-of-registers bisimulation
between the old and new source is required; equations (2) and (4)
replace that unnecessary demand.

\hypertarget{provenance-and-novelty}{%
\subsection{7. Provenance and novelty}\label{provenance-and-novelty}}

The original compact calculus and set theory are due to Mira-Cult, the
upstream verification to Jeff Ketchersid, and the machine compiler to
Andrew James Wade. ProveIt's pinned numerical calculus, enumeration,
safe-start and model adapters are inherited work.

Prior incoming material already uses rho6-derived ex falso and
quantified double-negation transport in
\nolinkurl{verification/double_b6_review.md} to compare a changed
quantifier row. The present contribution applies that method to the
\textbf{six exact fixed-bundle redexes}, enumerates all 64
original-layout masks, supplies their exact fine-token emission, and
explains the complete source/enumerator transfer needed to combine them
with a new binary optimization. It does not claim the underlying
double-negation principle as new.

Primary reference paths in the pinned ProveIt tree:

\begin{itemize}
\tightlist
\item
  \nolinkurl{Computability/BusyBeaver/Research/ZFC/candidate-routing245/README.md}
  --- literal baseline, count, seed, ordinary theorem and formal status.
\item
  \nolinkurl{Computability/BusyBeaver/Research/ZFC/native_theory254.md}
  --- exact numerical/native refutation comparison and full native
  formula schema.
\item
  \nolinkurl{Computability/BusyBeaver/Research/ZFC/dirty_program254.md}
  --- safe arbitrary control fields and structured source composition.
\item
  \nolinkurl{SetTheory/Cardinals/docs/reports/set-theory/zfc-closure-kernels/code/35-dirty-formal-src-BusyBeaver-ZFC-Code254.lean}
  --- numerical constructors, eight rows, exact bundle,
  \texttt{c254\_efq}.
\item
  The same source-35 prefix with \texttt{Lukasiewicz.lean} and
  \texttt{Kalmar.lean} --- implication derivations and
  deduction/completeness, without injectivity assumptions.
\item
  The same source-35 prefix with \texttt{Enum254.lean} and
  \texttt{EnumSafeStart254.lean} --- complete command contracts and
  cofinal padding.
\item
  \nolinkurl{docs/incoming/sub246-audit-and-alternatives-20261018.zip},
  \nolinkurl{research_sub246/verification/double_b6_review.md} and
  \nolinkurl{acceptance_checklist_and_transfer.md} --- prior quantified
  DNE argument and ordinary ZFC model interpretation audit.
\end{itemize}

The \nolinkurl{foundation_choice_route.md} note contains older
intermediate Lean status and must not by itself be read as the final
245-state ordinary calibration status. The baseline README and complete
incoming ordinary proof distinguish that status from the remaining exact
Lean endpoint.
